\documentclass[12pt, letterpaper]{article}

\usepackage[utf8]{inputenc}
\usepackage[T1]{fontenc}
\usepackage[spanish]{babel}
\usepackage{geometry}
\usepackage{amsmath, amssymb, amsthm}
\usepackage{booktabs}
\usepackage{array}
\usepackage{microtype}
\usepackage{setspace}
\usepackage{parskip}
\usepackage{titlesec}
\usepackage{fancyhdr}
\usepackage{enumitem}
\usepackage{bm}
\usepackage{mathtools}
\usepackage{tcolorbox}
\usepackage{hyperref}
\usepackage{pgfplots}
\usepackage{pgfplotstable}
\usepackage{xcolor}
\usepackage{colortbl}
\usepackage{multirow}
\usepackage{caption}

\pgfplotsset{compat=1.18}

\tcbuselibrary{skins, breakable}

\geometry{top=2.8cm, bottom=3.0cm, left=3.2cm, right=3.2cm}
\setstretch{1.20}

\pagestyle{fancy}
\fancyhf{}
\fancyhead[C]{\small\textsc{Econometr\'ia --- Gujarati \& Porter --- Ejercicio 1.1}}
\fancyfoot[C]{\small\thepage}
\renewcommand{\headrulewidth}{0.4pt}
\renewcommand{\footrulewidth}{0pt}

\titleformat{\section}
  {\large\bfseries\scshape}{\thesection.}{0.6em}{}
  [\vspace{-4pt}\rule{\linewidth}{0.4pt}]
\titleformat{\subsection}{\normalsize\bfseries}{\thesubsection.}{0.5em}{}
\titleformat{\subsubsection}{\normalsize\itshape}{}{0em}{}

\newtheorem{prop}{Proposici\'on}
\theoremstyle{remark}
\newtheorem{obs}{Observaci\'on}

\newtcolorbox{clave}{
  colback=white, colframe=black,
  boxrule=0.5pt, arc=3pt,
  left=8pt, right=8pt, top=6pt, bottom=6pt,
  breakable
}

\newtcolorbox{datebox}{
  colback=gray!4, colframe=black!30,
  boxrule=0.4pt, arc=2pt,
  left=8pt, right=8pt, top=5pt, bottom=5pt,
  breakable
}

% Colores para las series
\definecolor{col_canada}{HTML}{1b6ca8}
\definecolor{col_france}{HTML}{c0392b}
\definecolor{col_germany}{HTML}{27ae60}
\definecolor{col_italy}{HTML}{8e44ad}
\definecolor{col_japan}{HTML}{d35400}
\definecolor{col_uk}{HTML}{16a085}
\definecolor{col_usa}{HTML}{2c3e50}

\newcommand{\E}{\mathbb{E}}
\newcommand{\Var}{\operatorname{Var}}

\begin{document}

% ---------------------------------------------------------------
%  Portada
% ---------------------------------------------------------------
\begin{center}
  {\LARGE\bfseries Inflaci\'on en Siete Pa\'ises Industrializados}\\[6pt]
  {\LARGE\bfseries An\'alisis del \'Indice de Precios al Consumidor}\\[12pt]
  {\large\itshape C\'alculo de tasas, visualizaci\'on temporal}\\
  {\large\itshape y conclusiones comparativas (1974--1997)}\\[14pt]
  \rule{0.55\linewidth}{0.6pt}\\[8pt]
  {\normalsize Ejercicio~1.1 --- \textit{Econometr\'ia}, Gujarati \& Porter, 5.a ed.}\\[4pt]
  {\small\textsc{Fuente: Tabla~1.2 --- IPC base 1982--1984 $=$ 100}}
\end{center}

\vspace{10pt}

\noindent\textbf{El ejercicio.}\quad
La Tabla~1.2 de Gujarati \& Porter reporta el \'Indice de Precios al Consumidor
(IPC) de siete pa\'ises industrializados con base $1982\text{--}1984 = 100$.
Se solicita: \emph{(a)} calcular la tasa de inflaci\'on anual en cada pa\'is;
\emph{(b)} graficar las tasas de inflaci\'on en funci\'on del tiempo;
\emph{(c)} extraer conclusiones generales sobre la inflaci\'on en los siete pa\'ises;
y \emph{(d)} identificar el pa\'is con mayor variabilidad y explicarla.

% ---------------------------------------------------------------
\section{F\'ormula de c\'alculo}
% ---------------------------------------------------------------

La tasa de inflaci\'on anual se define como la variaci\'on porcentual del IPC
entre dos per\'iodos consecutivos:
\begin{equation}
  \pi_t
    = \frac{\text{IPC}_t - \text{IPC}_{t-1}}{\text{IPC}_{t-1}} \times 100
  \label{eq:inflacion}
\end{equation}
donde $\pi_t$ es la tasa de inflaci\'on en el a\~no $t$ (expresada en porcentaje)
e $\text{IPC}_t$ es el \'indice de precios al consumidor en ese mismo a\~no.
Esta f\'ormula es equivalente a la primera diferencia logar\'itmica para
variaciones peque\~nas:
\begin{equation}
  \pi_t \approx \bigl[\ln(\text{IPC}_t) - \ln(\text{IPC}_{t-1})\bigr]\times 100
  \label{eq:logdiff}
\end{equation}

% ---------------------------------------------------------------
\section{Datos del IPC y tasas de inflaci\'on calculadas}
% ---------------------------------------------------------------

\subsection{Tabla de IPC (Tabla~1.2 de Gujarati)}

\begin{datebox}
\small
\centering
\captionsetup{type=table}
\captionof{table}{\'Indice de Precios al Consumidor, base $1982\text{--}1984 = 100$}
\vspace{4pt}
\begin{tabular}{crrrrrrr}
\toprule
\textbf{A\~no} & \textbf{Canad\'a} & \textbf{Francia} & \textbf{Alemania} & \textbf{Italia} & \textbf{Jap\'on} & \textbf{R.U.} & \textbf{EE.UU.} \\
\midrule
1973 & 36.0 & 30.6 & 62.3 & 20.4 & 43.0 & 24.2 & 44.4 \\
1974 & 39.9 & 34.7 & 66.6 & 24.4 & 52.9 & 28.1 & 49.3 \\
1975 & 44.2 & 38.8 & 70.6 & 28.6 & 59.1 & 34.9 & 53.8 \\
1976 & 47.6 & 42.6 & 73.7 & 33.4 & 64.7 & 40.6 & 56.9 \\
1977 & 51.3 & 46.7 & 76.4 & 39.8 & 70.0 & 47.0 & 60.6 \\
1978 & 55.9 & 51.0 & 78.5 & 44.8 & 73.0 & 50.8 & 65.2 \\
1979 & 61.1 & 56.5 & 81.7 & 51.7 & 75.8 & 57.7 & 72.6 \\
1980 & 67.2 & 64.2 & 86.2 & 62.8 & 81.7 & 67.9 & 82.4 \\
1981 & 75.5 & 72.7 & 92.0 & 75.2 & 85.7 & 75.9 & 90.9 \\
1982 & 83.7 & 81.4 & 97.1 & 87.7 & 88.2 & 84.4 & 96.5 \\
1983 & 88.5 & 89.1 & 100.3 & 100.8 & 89.7 & 87.9 & 99.6 \\
1984 & 92.0 & 95.7 & 102.7 & 108.7 & 93.6 & 92.0 & 103.9 \\
1985 & 95.6 & 101.3 & 104.8 & 118.1 & 95.5 & 97.6 & 107.6 \\
1986 & 100.0 & 106.7 & 104.7 & 129.5 & 94.7 & 101.9 & 109.6 \\
1987 & 104.4 & 111.0 & 106.0 & 138.9 & 95.0 & 107.1 & 113.6 \\
1988 & 108.6 & 114.6 & 107.5 & 149.1 & 95.7 & 115.1 & 118.3 \\
1989 & 114.0 & 119.7 & 110.3 & 161.6 & 97.6 & 125.1 & 124.0 \\
1990 & 119.5 & 123.5 & 113.3 & 177.1 & 100.7 & 135.0 & 130.7 \\
1991 & 126.2 & 127.6 & 117.8 & 194.6 & 103.3 & 143.1 & 136.2 \\
1992 & 128.1 & 131.3 & 122.9 & 209.1 & 105.0 & 148.3 & 140.3 \\
1993 & 130.4 & 133.7 & 127.5 & 223.7 & 106.4 & 151.5 & 144.5 \\
1994 & 130.7 & 135.4 & 130.3 & 233.0 & 106.9 & 154.4 & 148.2 \\
1995 & 133.5 & 138.3 & 133.8 & 246.4 & 106.3 & 160.4 & 152.4 \\
1996 & 136.0 & 139.7 & 136.1 & 257.5 & 106.9 & 165.4 & 156.9 \\
1997 & 137.2 & 140.5 & 138.6 & 264.9 & 108.0 & 170.6 & 160.5 \\
\bottomrule
\end{tabular}
\end{datebox}

\subsection{Tasas de inflaci\'on calculadas (\%)}

Aplicando la ecuaci\'on~\eqref{eq:inflacion} a cada par de a\~nos consecutivos:

\begin{datebox}
\small
\centering
\captionsetup{type=table}
\captionof{table}{Tasas de inflaci\'on anual (\%), calculadas con \eqref{eq:inflacion}}
\vspace{4pt}
\begin{tabular}{crrrrrrrr}
\toprule
\textbf{A\~no} & \textbf{Canad\'a} & \textbf{Francia} & \textbf{Alemania} & \textbf{Italia} & \textbf{Jap\'on} & \textbf{R.U.} & \textbf{EE.UU.} \\
\midrule
1974 & 10.83 & 13.40 & 6.90 & 19.61 & 23.02 & 16.12 & 10.81 \\
1975 & 10.78 & 11.81 & 6.01 & 17.21 & 11.72 & 24.20 & 9.13 \\
1976 &  7.69 &  9.79 & 4.39 & 16.78 &  9.54 & 16.33 & 5.76 \\
1977 &  7.98 &  9.62 & 3.66 & 19.16 &  8.19 & 15.76 & 6.50 \\
1978 &  8.97 &  9.21 & 2.75 & 12.56 &  4.29 &  8.09 & 7.59 \\
1979 &  9.30 & 10.78 & 4.07 & 15.40 &  3.84 & 13.58 & 11.35 \\
1980 &  9.98 & 13.63 & 5.50 & 21.47 &  7.78 & 17.68 & 13.50 \\
1981 & 12.35 & 13.24 & 6.73 & 19.75 &  4.89 & 11.78 & 10.31 \\
1982 & 10.86 & 11.96 & 5.54 & 16.62 &  2.92 & 11.20 &  6.16 \\
1983 &  5.74 &  9.46 & 3.29 & 15.05 &  1.70 &  4.15 &  3.21 \\
1984 &  3.95 &  7.41 & 2.39 & 10.81 &  4.35 &  4.66 &  4.32 \\
1985 &  3.91 &  5.85 & 2.05 &  8.65 &  2.03 &  6.09 &  3.56 \\
1986 &  4.60 &  5.33 &-0.10 &  9.66 & -0.84 &  4.60 &  1.86 \\
1987 &  4.40 &  4.03 & 1.24 &  7.26 &  0.32 &  5.10 &  3.65 \\
1988 &  4.02 &  3.24 & 1.42 &  7.34 &  0.74 &  7.47 &  4.14 \\
1989 &  4.97 &  4.45 & 2.60 &  8.38 &  1.98 &  8.69 &  4.82 \\
1990 &  4.82 &  3.18 & 2.72 &  9.59 &  3.18 &  7.91 &  5.40 \\
1991 &  5.61 &  3.32 & 3.97 &  9.88 &  2.98 &  6.00 &  4.21 \\
1992 &  1.51 &  2.90 & 4.33 &  7.45 &  1.65 &  3.63 &  3.01 \\
1993 &  1.79 &  1.83 & 3.74 &  6.98 &  1.33 &  2.16 &  2.99 \\
1994 &  0.23 &  1.27 & 2.20 &  4.16 &  0.47 &  1.98 &  2.56 \\
1995 &  2.14 &  2.14 & 2.69 &  5.75 & -0.56 &  3.89 &  2.83 \\
1996 &  1.87 &  1.01 & 1.72 &  4.51 &  0.56 &  3.12 &  2.95 \\
1997 &  0.88 &  0.57 & 1.84 &  2.87 &  1.03 &  3.14 &  2.29 \\
\midrule
\textbf{Promedio} & 5.66 & 6.44 & 3.24 & 11.48 & 3.52 & 8.54 & 5.41 \\
\textbf{Desv.\ Est.} & 3.38 & 4.07 & 1.77 & 5.55 & 5.19 & 5.98 & 3.10 \\
\bottomrule
\end{tabular}
\end{datebox}

\medskip
\noindent\textit{Nota}: la inflaci\'on de 1974 se calcula respecto a 1973.

% ---------------------------------------------------------------
\section{Gr\'aficos de la tasa de inflaci\'on en funci\'on del tiempo}
% ---------------------------------------------------------------

Los siguientes gr\'aficos muestran la evoluci\'on de la tasa de inflaci\'on
en cada uno de los siete pa\'ises en el per\'iodo 1974--1997.

% --- Grafico 1: Canada, Francia, Alemania -----------------------
\begin{figure}[htbp]
\centering
\begin{tikzpicture}
\begin{axis}[
  width=13cm, height=6.5cm,
  xlabel={A\~no},
  ylabel={Tasa de inflaci\'on (\%)},
  xmin=1973, xmax=1998,
  ymin=-2, ymax=28,
  xtick={1974,1978,1982,1986,1990,1994,1997},
  xticklabel style={font=\small},
  yticklabel style={font=\small},
  grid=both,
  grid style={line width=0.2pt, draw=gray!30},
  major grid style={line width=0.4pt, draw=gray!50},
  legend pos=north east,
  legend style={font=\small, draw=black!30, fill=white, inner sep=4pt},
  tick align=outside,
  axis line style={line width=0.5pt},
  title style={font=\normalsize},
  title={Canad\'a, Francia y Alemania}
]
\addplot[color=col_canada, thick, mark=*, mark size=1.8pt] coordinates {
  (1974,10.83)(1975,10.78)(1976,7.69)(1977,7.98)(1978,8.97)
  (1979,9.30)(1980,9.98)(1981,12.35)(1982,10.86)(1983,5.74)
  (1984,3.95)(1985,3.91)(1986,4.60)(1987,4.40)(1988,4.02)
  (1989,4.97)(1990,4.82)(1991,5.61)(1992,1.51)(1993,1.79)
  (1994,0.23)(1995,2.14)(1996,1.87)(1997,0.88)
};
\addlegendentry{Canad\'a}
\addplot[color=col_france, thick, mark=square*, mark size=1.8pt] coordinates {
  (1974,13.40)(1975,11.81)(1976,9.79)(1977,9.62)(1978,9.21)
  (1979,10.78)(1980,13.63)(1981,13.24)(1982,11.96)(1983,9.46)
  (1984,7.41)(1985,5.85)(1986,5.33)(1987,4.03)(1988,3.24)
  (1989,4.45)(1990,3.18)(1991,3.32)(1992,2.90)(1993,1.83)
  (1994,1.27)(1995,2.14)(1996,1.01)(1997,0.57)
};
\addlegendentry{Francia}
\addplot[color=col_germany, thick, mark=triangle*, mark size=2pt] coordinates {
  (1974,6.90)(1975,6.01)(1976,4.39)(1977,3.66)(1978,2.75)
  (1979,4.07)(1980,5.50)(1981,6.73)(1982,5.54)(1983,3.29)
  (1984,2.39)(1985,2.05)(1986,-0.10)(1987,1.24)(1988,1.42)
  (1989,2.60)(1990,2.72)(1991,3.97)(1992,4.33)(1993,3.74)
  (1994,2.20)(1995,2.69)(1996,1.72)(1997,1.84)
};
\addlegendentry{Alemania}
\end{axis}
\end{tikzpicture}
\caption{Tasa de inflaci\'on anual: Canad\'a, Francia y Alemania (1974--1997).}
\end{figure}

% --- Grafico 2: Italia, Japon, RU, EEUU -------------------------
\begin{figure}[htbp]
\centering
\begin{tikzpicture}
\begin{axis}[
  width=13cm, height=6.5cm,
  xlabel={A\~no},
  ylabel={Tasa de inflaci\'on (\%)},
  xmin=1973, xmax=1998,
  ymin=-2, ymax=28,
  xtick={1974,1978,1982,1986,1990,1994,1997},
  xticklabel style={font=\small},
  yticklabel style={font=\small},
  grid=both,
  grid style={line width=0.2pt, draw=gray!30},
  major grid style={line width=0.4pt, draw=gray!50},
  legend pos=north east,
  legend style={font=\small, draw=black!30, fill=white, inner sep=4pt},
  tick align=outside,
  axis line style={line width=0.5pt},
  title style={font=\normalsize},
  title={Italia, Jap\'on, Reino Unido y EE.\,UU.}
]
\addplot[color=col_italy, thick, mark=diamond*, mark size=2pt] coordinates {
  (1974,19.61)(1975,17.21)(1976,16.78)(1977,19.16)(1978,12.56)
  (1979,15.40)(1980,21.47)(1981,19.75)(1982,16.62)(1983,15.05)
  (1984,10.81)(1985,8.65)(1986,9.66)(1987,7.26)(1988,7.34)
  (1989,8.38)(1990,9.59)(1991,9.88)(1992,7.45)(1993,6.98)
  (1994,4.16)(1995,5.75)(1996,4.51)(1997,2.87)
};
\addlegendentry{Italia}
\addplot[color=col_japan, thick, mark=pentagon*, mark size=2pt] coordinates {
  (1974,23.02)(1975,11.72)(1976,9.54)(1977,8.19)(1978,4.29)
  (1979,3.84)(1980,7.78)(1981,4.89)(1982,2.92)(1983,1.70)
  (1984,4.35)(1985,2.03)(1986,-0.84)(1987,0.32)(1988,0.74)
  (1989,1.98)(1990,3.18)(1991,2.98)(1992,1.65)(1993,1.33)
  (1994,0.47)(1995,-0.56)(1996,0.56)(1997,1.03)
};
\addlegendentry{Jap\'on}
\addplot[color=col_uk, thick, mark=o, mark size=2pt] coordinates {
  (1974,16.12)(1975,24.20)(1976,16.33)(1977,15.76)(1978,8.09)
  (1979,13.58)(1980,17.68)(1981,11.78)(1982,11.20)(1983,4.15)
  (1984,4.66)(1985,6.09)(1986,4.60)(1987,5.10)(1988,7.47)
  (1989,8.69)(1990,7.91)(1991,6.00)(1992,3.63)(1993,2.16)
  (1994,1.98)(1995,3.89)(1996,3.12)(1997,3.14)
};
\addlegendentry{Reino Unido}
\addplot[color=col_usa, thick, dashed, mark=x, mark size=2.5pt, mark options={solid}] coordinates {
  (1974,10.81)(1975,9.13)(1976,5.76)(1977,6.50)(1978,7.59)
  (1979,11.35)(1980,13.50)(1981,10.31)(1982,6.16)(1983,3.21)
  (1984,4.32)(1985,3.56)(1986,1.86)(1987,3.65)(1988,4.14)
  (1989,4.82)(1990,5.40)(1991,4.21)(1992,3.01)(1993,2.99)
  (1994,2.56)(1995,2.83)(1996,2.95)(1997,2.29)
};
\addlegendentry{EE.\,UU.}
\end{axis}
\end{tikzpicture}
\caption{Tasa de inflaci\'on anual: Italia, Jap\'on, Reino Unido y EE.\,UU.\ (1974--1997).}
\end{figure}

% --- Grafico 3: Comparativo todos -------------------------------
\begin{figure}[htbp]
\centering
\begin{tikzpicture}
\begin{axis}[
  width=14cm, height=7cm,
  xlabel={A\~no},
  ylabel={Tasa de inflaci\'on (\%)},
  xmin=1973, xmax=1998,
  ymin=-2, ymax=28,
  xtick={1974,1976,1978,1980,1982,1984,1986,1988,1990,1992,1994,1996},
  xticklabel style={font=\scriptsize, rotate=45, anchor=east},
  yticklabel style={font=\small},
  grid=both,
  grid style={line width=0.2pt, draw=gray!25},
  major grid style={line width=0.35pt, draw=gray!45},
  legend columns=4,
  legend style={
    font=\footnotesize, draw=black!30, fill=white,
    inner sep=4pt, at={(0.5,-0.22)}, anchor=north
  },
  tick align=outside,
  axis line style={line width=0.5pt},
  title style={font=\normalsize},
  title={Comparativo: los siete pa\'ises (1974--1997)}
]
\addplot[color=col_canada,  thick, mark=*, mark size=1.5pt] coordinates {
  (1974,10.83)(1975,10.78)(1976,7.69)(1977,7.98)(1978,8.97)(1979,9.30)
  (1980,9.98)(1981,12.35)(1982,10.86)(1983,5.74)(1984,3.95)(1985,3.91)
  (1986,4.60)(1987,4.40)(1988,4.02)(1989,4.97)(1990,4.82)(1991,5.61)
  (1992,1.51)(1993,1.79)(1994,0.23)(1995,2.14)(1996,1.87)(1997,0.88)
};
\addlegendentry{Canad\'a}
\addplot[color=col_france, thick, mark=square*, mark size=1.5pt] coordinates {
  (1974,13.40)(1975,11.81)(1976,9.79)(1977,9.62)(1978,9.21)(1979,10.78)
  (1980,13.63)(1981,13.24)(1982,11.96)(1983,9.46)(1984,7.41)(1985,5.85)
  (1986,5.33)(1987,4.03)(1988,3.24)(1989,4.45)(1990,3.18)(1991,3.32)
  (1992,2.90)(1993,1.83)(1994,1.27)(1995,2.14)(1996,1.01)(1997,0.57)
};
\addlegendentry{Francia}
\addplot[color=col_germany, thick, mark=triangle*, mark size=1.8pt] coordinates {
  (1974,6.90)(1975,6.01)(1976,4.39)(1977,3.66)(1978,2.75)(1979,4.07)
  (1980,5.50)(1981,6.73)(1982,5.54)(1983,3.29)(1984,2.39)(1985,2.05)
  (1986,-0.10)(1987,1.24)(1988,1.42)(1989,2.60)(1990,2.72)(1991,3.97)
  (1992,4.33)(1993,3.74)(1994,2.20)(1995,2.69)(1996,1.72)(1997,1.84)
};
\addlegendentry{Alemania}
\addplot[color=col_italy, thick, mark=diamond*, mark size=1.8pt] coordinates {
  (1974,19.61)(1975,17.21)(1976,16.78)(1977,19.16)(1978,12.56)(1979,15.40)
  (1980,21.47)(1981,19.75)(1982,16.62)(1983,15.05)(1984,10.81)(1985,8.65)
  (1986,9.66)(1987,7.26)(1988,7.34)(1989,8.38)(1990,9.59)(1991,9.88)
  (1992,7.45)(1993,6.98)(1994,4.16)(1995,5.75)(1996,4.51)(1997,2.87)
};
\addlegendentry{Italia}
\addplot[color=col_japan, thick, mark=pentagon*, mark size=1.8pt] coordinates {
  (1974,23.02)(1975,11.72)(1976,9.54)(1977,8.19)(1978,4.29)(1979,3.84)
  (1980,7.78)(1981,4.89)(1982,2.92)(1983,1.70)(1984,4.35)(1985,2.03)
  (1986,-0.84)(1987,0.32)(1988,0.74)(1989,1.98)(1990,3.18)(1991,2.98)
  (1992,1.65)(1993,1.33)(1994,0.47)(1995,-0.56)(1996,0.56)(1997,1.03)
};
\addlegendentry{Jap\'on}
\addplot[color=col_uk, thick, mark=o, mark size=1.8pt] coordinates {
  (1974,16.12)(1975,24.20)(1976,16.33)(1977,15.76)(1978,8.09)(1979,13.58)
  (1980,17.68)(1981,11.78)(1982,11.20)(1983,4.15)(1984,4.66)(1985,6.09)
  (1986,4.60)(1987,5.10)(1988,7.47)(1989,8.69)(1990,7.91)(1991,6.00)
  (1992,3.63)(1993,2.16)(1994,1.98)(1995,3.89)(1996,3.12)(1997,3.14)
};
\addlegendentry{Reino Unido}
\addplot[color=col_usa, thick, dashed, mark=x, mark size=2pt, mark options={solid}] coordinates {
  (1974,10.81)(1975,9.13)(1976,5.76)(1977,6.50)(1978,7.59)(1979,11.35)
  (1980,13.50)(1981,10.31)(1982,6.16)(1983,3.21)(1984,4.32)(1985,3.56)
  (1986,1.86)(1987,3.65)(1988,4.14)(1989,4.82)(1990,5.40)(1991,4.21)
  (1992,3.01)(1993,2.99)(1994,2.56)(1995,2.83)(1996,2.95)(1997,2.29)
};
\addlegendentry{EE.\,UU.}
\end{axis}
\end{tikzpicture}
\caption{Tasas de inflaci\'on anual en los siete pa\'ises industrializados, 1974--1997.
  Se observan con claridad los dos picos de los choques petroleros y la convergencia
  hacia tasas bajas durante la d\'ecada de 1990.}
\end{figure}

% ---------------------------------------------------------------
\section{Conclusiones generales}
% ---------------------------------------------------------------

\begin{obs}
  \textbf{Dos episodios de inflaci\'on elevada.} Los gr\'aficos muestran dos
  picos claramente identificables en todos los pa\'ises: 1974--1975 y
  1979--1981. Estos episodios coinciden con los choques de precios del
  petr\'oleo provocados por la OPEP en 1973 y 1979, respectivamente.
  El primer choque elev\'o las tasas de inflaci\'on de los pa\'ises
  importadores de petr\'oleo de forma sincronizada y abrupta.
\end{obs}

\begin{obs}
  \textbf{Tendencia generalizada a la baja desde 1982.} A partir de 1982,
  las tasas de inflaci\'on de los siete pa\'ises iniciaron una reducci\'on
  sostenida. Este descenso es atribuible a la adopci\'on de pol\'iticas
  monetarias restrictivas (especialmente la elecci\'on del control de la
  inflaci\'on como objetivo prioritario de los bancos centrales) y, en
  particular, a las pol\'iticas del presidente de la Reserva Federal Paul
  Volcker en EE.\,UU., que sirvieron de referencia internacional.
\end{obs}

\begin{obs}
  \textbf{Convergencia en la d\'ecada de 1990.} Para mediados de los a\~nos
  noventa, las tasas de inflaci\'on de los siete pa\'ises convergieron hacia
  un rango de $1\%$ a $4\%$ anual. Esta convergencia refleja tanto la
  globalizaci\'on de las pol\'iticas de estabilizaci\'on como la mayor
  credibilidad de los bancos centrales independientes.
\end{obs}

\begin{obs}
  \textbf{Heterogeneidad en los niveles.} A pesar de la tendencia com\'un,
  existen diferencias sustanciales en los \emph{niveles} de inflaci\'on a lo
  largo de todo el per\'iodo. Italia mantuvo tasas sistem\'aticamente m\'as
  altas que el resto; Alemania y Jap\'on presentaron las tasas m\'as bajas.
  Esta heterogeneidad refleja diferencias estructurales en las instituciones
  monetarias, la rigidez de los mercados laborales y el grado de apertura
  comercial de cada econom\'ia.
\end{obs}

% ---------------------------------------------------------------
\section{Pa\'is con la tasa de inflaci\'on m\'as variable}
% ---------------------------------------------------------------

\subsection{Criterios de variabilidad}

Para cuantificar la variabilidad de la tasa de inflaci\'on se utilizan dos
estad\'isticos. La \textbf{desviaci\'on est\'andar} mide la dispersi\'on
absoluta alrededor de la media:
\begin{equation}
  s_\pi = \sqrt{\frac{1}{T-1}\sum_{t=1}^{T}(\pi_t - \bar\pi)^2}
  \label{eq:std}
\end{equation}
El \textbf{coeficiente de variaci\'on} mide la dispersi\'on relativa y
permite comparar pa\'ises con promedios muy distintos:
\begin{equation}
  \text{CV} = \frac{s_\pi}{\bar\pi} \times 100
  \label{eq:cv}
\end{equation}

\subsection{Resultados}

\begin{center}
\begin{tabular}{lrrrr}
\toprule
\textbf{Pa\'is} & \textbf{Promedio} (\%) & \textbf{M\'aximo} (\%) & $s_\pi$ & \textbf{CV} (\%) \\
\midrule
Canad\'a   &  5.66 & 12.35 & 3.38 &  59.7 \\
Francia    &  6.44 & 13.63 & 4.07 &  63.2 \\
Alemania   &  3.24 &  6.90 & 1.77 &  54.6 \\
Italia     & 11.48 & 21.47 & 5.55 &  48.3 \\
\textbf{Jap\'on}  &  \textbf{3.52} & \textbf{23.02} & \textbf{5.19} & \textbf{147.4} \\
\textbf{R.U.}     &  \textbf{8.54} & \textbf{24.20} & \textbf{5.98} &  \textbf{70.0} \\
EE.\,UU.   &  5.41 & 13.50 & 3.10 &  57.3 \\
\bottomrule
\end{tabular}
\end{center}

\subsection{Interpretaci\'on}

\begin{clave}
\textbf{El pa\'is con la tasa de inflaci\'on m\'as variable es Jap\'on},
medido por el coeficiente de variaci\'on ($\text{CV} = 147.4\%$), y
\textbf{el Reino Unido} si se mide por la desviaci\'on est\'andar absoluta
($s_\pi = 5.98$).

\medskip
\textbf{Explicaci\'on para Jap\'on.} La extrema variabilidad japonesa tiene
una causa precisa: el pico de $23.02\%$ en 1974, el m\'as alto de los siete
pa\'ises en ese a\~no, seguido de una ca\'ida vertiginosa hasta $1.70\%$ en
1983 y casi cero en los noventa. Jap\'on era en 1973 la econom\'ia
desarrollada m\'as dependiente de las importaciones de petr\'oleo (m\'as del
$70\%$ de su energ\'ia primaria proced\'ia del exterior). El choque de la
OPEP impact\'o su estructura de costos de forma desproporcionada respecto a
los dem\'as pa\'ises. Sin embargo, la respuesta de pol\'itica monetaria del
Banco de Jap\'on fue tambi\'en la m\'as agresiva y exitosa: en menos de una
d\'ecada, la econom\'ia japonesa pas\'o de la inflaci\'on m\'as alta a la
m\'as baja del grupo.

\medskip
\textbf{Explicaci\'on para el Reino Unido.} El Reino Unido tuvo el mayor
pico de inflaci\'on absoluto: $24.20\%$ en 1975, un a\~no despu\'es del
primer choque petrolero. Esto se explica por la combinaci\'on de la
dependencia energ\'etica con la debilidad de la libra esterlina en ese
per\'iodo y con acuerdos salariales indexados que amplificaron el efecto
del choque de oferta (\textit{segunda vuelta de la inflaci\'on}). La inflaci\'on
brit\'anica se redujo gradualmente bajo las pol\'iticas de austeridad del
gobierno de Margaret Thatcher a partir de 1979.
\end{clave}

\appendix

\section{C\'alculo detallado: ejemplo para Jap\'on (1974--1976)}

Aplicando la ecuaci\'on~\eqref{eq:inflacion}:
\begin{align}
  \pi_{1974}^{\text{JP}}
    &= \frac{52.9 - 43.0}{43.0}\times 100
     = \frac{9.9}{43.0}\times 100
     = 23.02\% \\[4pt]
  \pi_{1975}^{\text{JP}}
    &= \frac{59.1 - 52.9}{52.9}\times 100
     = \frac{6.2}{52.9}\times 100
     = 11.72\% \\[4pt]
  \pi_{1976}^{\text{JP}}
    &= \frac{64.7 - 59.1}{59.1}\times 100
     = \frac{5.6}{59.1}\times 100
     = 9.47\%
\end{align}
La ca\'ida de $23.02\%$ a $9.47\%$ en dos a\~nos ilustra la rapidez con
que la pol\'itica monetaria japonesa contuvo el choque inflacionario.

\section{Nota sobre la diferencia logar\'itmica}

Para valores peque\~nos de $\pi_t$, la aproximaci\'on logar\'itmica
\eqref{eq:logdiff} es muy precisa. Para el caso de Jap\'on en 1974:
\begin{equation}
  \ln(52.9) - \ln(43.0) = 3.9679 - 3.7612 = 0.2067 \approx 20.67\%
\end{equation}
La diferencia con el valor exacto ($23.02\%$) es de $2.35$ puntos
porcentuales, lo que muestra que la aproximaci\'on logar\'itmica
\emph{subestima} la inflaci\'on real cuando las tasas son elevadas.
Para an\'alisis de series de tiempo de largo plazo se prefiere la
f\'ormula exacta \eqref{eq:inflacion}.

\vspace{14pt}
\noindent\rule{\linewidth}{0.4pt}

\noindent{\small\textit{Nota.} Los datos del IPC provienen de la Tabla~1.2
de Gujarati \& Porter (5.a ed.). Las tasas de inflaci\'on calculadas en la
Tabla~2 de este documento pueden diferir levemente de otras fuentes seg\'un
el criterio de redondeo del IPC original. El c\'odigo de c\'alculo en \textsf{R}
est\'a disponible en \texttt{notebooks/01\_naturaleza\_regresion/}.}

\end{document}
